\documentclass[11pt]{article}
\usepackage[T1]{fontenc}
\usepackage[margin=1in]{geometry}
\usepackage{amsmath,amssymb,amsthm}
\usepackage{microtype}
\usepackage[hidelinks]{hyperref}
\newtheorem{theorem}{Theorem}
\newcommand{\R}{\mathbb{R}}
\newcommand{\Ptwo}{\mathcal P_2(\R^n)}
\newcommand{\Id}{\mathrm{Id}}
\newcommand{\Law}{\operatorname{Law}}
\begin{document}

\begin{center}
{\Large\bfseries A Convex-Order Characterization of All Dual Probabilistic Frames}\par
\vspace{0.5em}
{\large A literature-implied full answer to the question in arXiv:2505.13885}\par
\vspace{0.4em}
Run \texttt{fa\_banach\_001} --- lane 17 \qquad 13 August 2026
\end{center}

\medskip
\noindent\fbox{\begin{minipage}{0.94\linewidth}
\textbf{Status.} Full characterization, obtained by combining conditional
barycenters with Strassen's martingale theorem. A measure is dual to a fixed
probabilistic frame precisely when it is a convex-order mean-preserving spread
of an admissible pushforward-type dual.
\end{minipage}}

\section{The source question}

Chen, King, and Shonkwiler's \emph{Approximately Dual and Pseudo-Dual
Probabilistic Frames} \cite{CKS} defines a dual of a probabilistic frame
$\mu\in\Ptwo$ to be a measure $\nu\in\Ptwo$ for which some coupling
$\gamma\in\Gamma(\mu,\nu)$ satisfies
\begin{equation}\label{eq:dual}
 \int_{\R^n\times\R^n}xy^t\,d\gamma(x,y)=\Id.
\end{equation}
The Introduction (printed and PDF page~3) states that characterizing all such
$\nu$ for a given $\mu$ remains open. The statement occurs in the paragraph beginning
``Characterizing all the dual probabilistic frames''; the same challenge is
repeated at the beginning of the section on approximately dual frames.

\section{Exact characterization}

Let
\[
 S_\mu:=\int_{\R^n}xx^t\,d\mu(x)
\]
be the positive definite frame operator, and define
\[
 \mathcal H_\mu:=\left\{h\in L^2(\mu;\R^n):
       \int_{\R^n}x h(x)^t\,d\mu(x)=0\right\}.
\]
For probability measures $\eta,\nu$ with finite first moments, write
$\eta\preceq_{\rm cx}\nu$ when
$\int\varphi\,d\eta\leq\int\varphi\,d\nu$ for every convex $\varphi$
for which the integrals are defined.

\begin{theorem}\label{thm:classification}
For every probabilistic frame $\mu$ on $\R^n$, its complete class of dual
probabilistic frames is
\begin{equation}\label{eq:class}
 \mathcal D(\mu)=
 \bigcup_{h\in\mathcal H_\mu}
 \left\{\nu\in\Ptwo:
   \big(S_\mu^{-1}\Id+h\big)_\#\mu\preceq_{\rm cx}\nu\right\}.
\end{equation}
More explicitly, $\nu$ is dual to $\mu$ if and only if there are
$h\in\mathcal H_\mu$ and a probability kernel $K_z(dy)$ such that, with
$T(x)=S_\mu^{-1}x+h(x)$,
\[
 \int y\,K_z(dy)=z
 \quad\text{for }T_\#\mu\text{-a.e. }z,
 \qquad
 \nu(B)=\int K_{T(x)}(B)\,d\mu(x).
\]
Thus every dual is a mean-preserving spread of a pushforward-type dual, and
every such spread is a dual.
\end{theorem}

\begin{proof}
Suppose first that $\nu$ is dual to $\mu$, witnessed by $\gamma$. Disintegrate
\[
 d\gamma(x,y)=d\mu(x)\,K_x(dy)
\]
and let $T(x):=\int y\,K_x(dy)$. Conditional Jensen and $\nu\in\Ptwo$ give
$T\in L^2(\mu;\R^n)$ and
\[
 T_\#\mu\preceq_{\rm cx}\nu.
\]
Moreover, \eqref{eq:dual} becomes
\[
 \int xT(x)^t\,d\mu(x)=\Id.
\]
Put $h(x)=T(x)-S_\mu^{-1}x$. Since $S_\mu$ is symmetric,
\[
 \int xh(x)^t\,d\mu(x)
 =\Id-S_\mu(S_\mu^{-1})^t=0.
\]
Hence $h\in\mathcal H_\mu$, proving the inclusion from left to right in
\eqref{eq:class}.

Conversely, take $h\in\mathcal H_\mu$, put
$T(x)=S_\mu^{-1}x+h(x)$ and $\eta=T_\#\mu$, and suppose
$\eta\preceq_{\rm cx}\nu$. Strassen's martingale theorem, in the explicit
multidimensional kernel form \cite[Theorem~2.3]{BB}, supplies a probability
kernel $K_z(dy)$ with first
marginal $\eta$, second marginal $\nu$, and barycenter
$\int yK_z(dy)=z$. Define
\[
 d\gamma(x,y):=d\mu(x)\,K_{T(x)}(dy).
\]
Then $\gamma\in\Gamma(\mu,\nu)$ and
\begin{align*}
 \int xy^t\,d\gamma(x,y)
 &=\int x\left(\int y\,K_{T(x)}(dy)\right)^t d\mu(x)\\
 &=\int xT(x)^t\,d\mu(x)\\
 &=S_\mu(S_\mu^{-1})^t+
   \int xh(x)^t\,d\mu(x)=\Id.
\end{align*}
Thus $\nu$ is a dual. (The same Cauchy--Schwarz argument used in the source
then shows automatically that $\nu$ is a probabilistic frame.)
\end{proof}

The point-mass case checks the orientation. If $\mu=\delta_a$ on $\R$ with
$a\ne0$, then $\mathcal H_\mu=\{0\}$ and the theorem says
$\nu$ is dual exactly when $\delta_{1/a}\preceq_{\rm cx}\nu$, equivalently
when $\nu$ has mean $1/a$. This recovers the complete point-mass description
in Chen and Schmoll \cite[Proposition~5.12]{CS}.

\section{Provenance, search, and scope}

This is a literature-implied answer. Strassen's 1965 theorem says that convex
order is equivalent to the existence of a martingale coupling; it does not
mention frames. The theorem is reproduced in explicit multidimensional kernel
form as Theorem~2.3 in \cite{BB}. The source paper gives the pushforward
condition (its Lemma~3.7), but does not pass from arbitrary couplings to
conditional barycenters and convex order. The combination above is the direct
identification that closes the stated characterization problem.

A bounded search through 13 August 2026 checked the full texts of
arXiv:2505.13885 and arXiv:2501.02602 for the terms
``convex order'', ``martingale'', ``Strassen'', ``barycenter'', and
``conditional expectation'', and searched the exact phrases
``dual probabilistic frame convex order'', ``transport dual martingale'', and
close variants. No prior statement of Theorem~\ref{thm:classification} was
found. The source authors do not cite Strassen here and did not explicitly
know they were posing an instance of his theorem. No priority claim is made.

The theorem completely characterizes which second marginal measures are
duals, in every finite dimension. It does not classify extreme points of this
nonclosed class, choose a canonical admissible $h$ for a prescribed $\nu$, or
provide a numerical martingale-transport algorithm.

\begin{sloppypar}
\begin{thebibliography}{9}
\raggedright
\bibitem{CKS}
D. Chen, E. J. King, and C. Shonkwiler,
\emph{Approximately Dual and Pseudo-Dual Probabilistic Frames},
arXiv:2505.13885 (2025). See the Introduction, p.~3, and Lemma~3.7, p.~9.

\bibitem{CS}
D. Chen and M. Schmoll,
\emph{Probabilistic Frames and Wasserstein Distances},
arXiv:2501.02602v1 (2025). See Section~5 and Proposition~5.12.

\bibitem{Strassen}
V. Strassen,
\emph{The Existence of Probability Measures with Given Marginals},
Ann. Math. Statist. \textbf{36} (1965), 423--439.
See Theorem~8.

\bibitem{BB}
K. Bo{\l}botowski and G. Bouchitt\'e,
\emph{Kantorovich--Rubinstein Duality Theory for the Hessian},
arXiv:2412.00516 (2024). See Theorem~2.3, printed and PDF page~15.
\end{thebibliography}
\end{sloppypar}

\end{document}
